\documentclass[11pt]{article}
\usepackage[margin=1in]{geometry}
\usepackage{amsmath,amssymb,amsthm,booktabs,microtype}
\usepackage[colorlinks=true,linkcolor=blue,urlcolor=blue]{hyperref}
\newtheorem{theorem}{Theorem}
\newtheorem{lemma}{Lemma}
\newcommand{\norm}[1]{\left\lVert #1\right\rVert_\infty}
\newcommand{\dist}{\operatorname{dist}}
\newcommand{\R}{\mathbb R}
\title{A self-similar Lorentz counterexample:\\an optimized proof and exact certificate}
\author{}
\date{October 28, 2026}
\begin{document}
\maketitle
\begin{abstract}
We construct a continuous nonodd function on $[-1,1]$ whose best uniform
polynomial approximant of every degree vanishes at the origin. A finite
certificate handles degrees zero through six; cubic self-similarity handles
all remaining degrees. Keeping positive and negative contact sets separate
reduces the analytic persistence argument to one inequality and one smallness
condition. The construction works for every $0<\varepsilon\le 4\cdot10^{-8}$,
including the earlier choice $10^{-10}$. Exact finite checks are reproducible
from the accompanying Lean coefficient tables.
\end{abstract}

\section{Construction and reduction to seven degrees}
Write $I=[-1,1]$, let $T_k(\cos\theta)=\cos(k\theta)$, and let $P_n f$ denote
the unique best uniform approximant to $f\in C(I)$ among polynomials of degree
at most $n$. We address the origin interpolation implication discussed
in~\cite{sv}. Set $z=2\cos(\pi/36)$ and
\begin{align*}
107\beta={}&-60+82z+1086z^2+1576z^3-2932z^4-3433z^5\\
 &+2002z^6+2243z^7-510z^8-587z^9+44z^{10}+54z^{11}.
\end{align*}
In particular $419/1000<\beta<420/1000$. Define
\[
 h=T_4-1-2T_5+\beta T_7,\qquad q_\varepsilon=T_3+\varepsilon h,
 \qquad f_\varepsilon=\sum_{j\ge0}(-1/3)^j q_\varepsilon\circ T_{3^j}.
\]
The series converges uniformly on $I$, so $f_\varepsilon$ is continuous.

\begin{theorem}\label{thm:main}
For every $0<\varepsilon\le 1/25000000$, the above function satisfies
\[
 (P_n f_\varepsilon)(0)=0\quad(n\ge0),\qquad
 f_\varepsilon(1/\sqrt2)+f_\varepsilon(-1/\sqrt2)=-3\varepsilon.
\]
Thus it is not odd.
\end{theorem}

The reduction to finitely many degrees uses the identity
\begin{equation}\label{eq:composition}
 P_n(g\circ T_M)=(P_{\lfloor n/M\rfloor}g)\circ T_M\qquad(M\ge1).
\end{equation}
To prove it, average a competitor $p(\cos\theta)$ over
$\theta+2\pi j/M$, $0\le j<M$. The target is invariant under these shifts,
and the average cannot increase its uniform error. Only cosine frequencies
divisible by $M$ remain, giving $r(\cos M\theta)$ with
$\deg r\le\lfloor n/M\rfloor$. Surjectivity of $T_M:I\to I$ gives equality
of the optimal errors, and uniqueness of best approximation gives
\eqref{eq:composition}.

The functional equation and its consequence for $n\ge7$ are
\begin{equation}\label{eq:induction}
 f_\varepsilon=q_\varepsilon-\tfrac13 f_\varepsilon\circ T_3,
 \qquad
 P_n f_\varepsilon=q_\varepsilon-\tfrac13
 (P_{\lfloor n/3\rfloor}f_\varepsilon)\circ T_3.
\end{equation}
Here we used polynomial translation, scalar homogeneity, and
$\deg q_\varepsilon\le7$. Since $q_\varepsilon(0)=T_3(0)=0$, induction reduces
the origin condition to $0\le n\le6$.

\section{Baseline residuals and the finite certificate}
Put
\[
 a=1/\sqrt2,\quad
 G=\sum_{j\ge0}(-1/3)^j T_{3^j},\quad F=G\circ T_3,\quad
 H=\sum_{j\ge0}(-1/3)^j h\circ T_{3^j}.
\]
Then $f_\varepsilon=F+\varepsilon H$. Pairing consecutive terms gives
\[
 G=\sum_{j\ge0}9^{-j}g\circ T_{9^j},\qquad
 g(x)=2x-\tfrac43x^3.
\]
The unique maximum of $g$ on $I$ is at $a$, its unique minimum is at $-a$,
and $T_9(\pm a)=\pm a$. Hence $G$ has the same unique extrema and
\[
 E:=\norm G=\frac{g(a)}{1-1/9}=\frac{3a}{2}.
\]
Oddness gives $P_0G=0$. From $G=x-G\circ T_3/3$ and
\eqref{eq:composition}, $P_1G=P_2G=x$. Thus for the seven base degrees,
\begin{equation}\label{eq:baseline}
 p_n^0=P_nF=\begin{cases}0&n\le2,\\ T_3&3\le n\le6,\end{cases}
 \qquad b_n:=F-p_n^0=
 \begin{cases}G\circ T_3&n\le2,\\-G\circ T_9/3&3\le n\le6.\end{cases}
\end{equation}
Their norms $E_n$ are respectively $E$ and $E/3$. Their contact sets $S_n$
are respectively the zeros of $T_6$ and $T_{18}$.

Here are exact definitions of the finite certificate. Use the increasing nodes
\[
 X_j=-\cos\frac{(2j-1)\pi}{36}\quad(1\le j\le18),\qquad
 (C_1,\ldots,C_6)=(X_2,X_5,X_8,X_{11},X_{14},X_{17}).
\]
The residual signs on the fine grid repeat $(+,-,-,+)$; call them $\sigma_j$.
On the coarse grid they repeat $(-,+,+,-)$; call them $\tau_j$.
The two-cycle $T_3(a)=-a$, $T_3(-a)=a$ gives
\[
 H(\pm a)=-\tfrac32\pm\tfrac32(2+\beta)a.
\]
Iterating $H=h-H\circ T_3/3$ twice, and using $T_{21}=-T_{15}$ at the
zeros of $T_{18}$, yields $H(X_j)=R_\beta(X_j)$, where
\begin{equation}\label{eq:gridvalue}
 R_\beta=T_4-\tfrac13T_{12}-\tfrac56-2T_5+\beta T_7
                +\tfrac{2+\beta}{3}T_{15}+\tfrac{2+\beta}{6}T_9.
\end{equation}

Define corrections $v_n$ and slopes $\lambda_n$ as follows:
\begin{align*}
 v_0&=0,&\lambda_0&=\tfrac{\sqrt6}{4}(2+\beta),\\
 v_1=v_2&=-(2+\beta)x,&\lambda_1=\lambda_2&=\tfrac{2+\beta}{2\sqrt2},\\
 v_3&=A_1x+A_2x^2+A_3x^3,\\
 v_4&=T_4-1+B_1T_1+B_3T_3,\\
 v_5=v_6&=T_4-1+D_1T_1+D_3T_3+D_5T_5,&\lambda_6&=\lambda_5.
\end{align*}
The remaining coefficients and slopes are the solutions of
\begin{equation}\label{eq:systems}
 v_n(X_j)+\sigma_j\lambda_n=R_\beta(X_j),\qquad
 \begin{array}{c|l}
 n&j\\\hline
 3&1,3,8,10\\
 4&10,16,18\\
 5&10,13,15,16.
 \end{array}
\end{equation}
The three defining matrices are nonsingular, as certified below.
In particular $\deg v_n\le n$ and $v_n(0)=0$.

For a contact $\xi\in S_n$, write $s(\xi)=\operatorname{sgn} b_n(\xi)$.
The exact finite calculation establishes
\begin{equation}\label{eq:finite}
 s(\xi)\bigl(H(\xi)-v_n(\xi)\bigr)\le\lambda_n,
 \qquad 0<\lambda_n<2,
\end{equation}
with the following complete zero-slack sets. Indices refer to the indicated grid.
\begin{center}
\begin{tabular}{cclc}
\toprule
$n$&grid&zero-slack indices&other slacks exceed\\\midrule
0&$C$&$3,4$&$127/50$\\
1,2&$C$&$1,3,4,6$&$363/100$\\
3&$X$&$1,3,8,10,16$&$9/100$\\
4&$X$&$1,3,8,9,10,16,18$&$7/10$\\
5,6&$X$&$3,4,6,9,10,13,15,16$&$1/5$\\\bottomrule
\end{tabular}
\end{center}
For $n=0,2,3,6$, use the whole listed set. For $n=1,4,5$ use respectively
$\{1,3,4\}$, $\{1,3,9,10,16,18\}$, and $\{3,4,6,9,10,13,15\}$.
Each selected set contains $n+2$ increasing nodes with alternating signs.

The sums of the absolute Chebyshev coefficients of $v_n$ are respectively
\[
 0,\quad <5/2,\quad <5/2,\quad <9/4,\quad <13/4,\quad <11/2,\quad <11/2.
\]
Consequently $\norm{v_n}<6$ and $\norm{v_n'}<216$, since
$\norm{T_k'}=k^2$ and $\deg v_n\le6$.

\subsection*{Exact verification of the finite assertions}
All the numbers above lie in $\mathbb Q[z]$, with arithmetic reduced modulo
\[
 m(z)=z^{12}-12z^{10}+54z^8-112z^6+105z^4-36z^2+1
      =(z^6-6z^4+9z^2-2)^2-3.
\]
The last identity and $z=2\cos(\pi/36)$ give $m(z)=0$. There is a unique root
in $(1.98,2)$: the inner polynomial is positive there and has derivative
$6z(z^2-1)(z^2-3)>0$. Rational endpoint evaluations place the root in
\[
 [199238939/10^8,\;199238940/10^8].
\]
The accompanying \texttt{scripts/check\_certificate.py} reads exact rational
coefficient lists from \texttt{CertificateData.lean} and
\texttt{CertificateGrid.lean}. It reconstructs the nodes by the Chebyshev
recurrence, verifies \eqref{eq:gridvalue}, the functional equation on the grid,
every slack identity, every equality set, alternation, coefficient norms, and
the nonzero determinants in \eqref{eq:systems}. Equality is tested as equality
of reduced rational polynomials. One hundred rational bisections, followed by
rational interval Horner evaluation, certify the strict inequalities. No
floating-point tolerance is used. This also verifies that the stored
coefficients are the uniquely defined solutions of \eqref{eq:systems}.

\section{Signed cusps and a global bound relative to contacts}
The two estimates needed to turn \eqref{eq:finite} into uniform approximation
are proved here with explicit constants.

\begin{lemma}\label{lem:gap}
For $s\in\{-1,1\}$ set $S_{n,s}=\{\xi\in I:b_n(\xi)=sE_n\}$.
For $0\le n\le6$ and $x\in I$,
\[
 E_n-sb_n(x)\ge\frac1{250}\dist(x,S_{n,s}).
\]
\end{lemma}
\begin{proof}
First we prove $E-sG(y)\ge c_G|y-sa|$ on $I$, where
$c_G=\pi(\sqrt2-1)/324$. The factorization
\[
 g(a)-g(y)=\tfrac43(y-a)^2(y+\sqrt2)
       \ge C_0(y-a)^2,\qquad C_0=\tfrac43(\sqrt2-1),
\]
makes every term of $E-G(y)=\sum_{j\ge0}9^{-j}(g(a)-g(T_{9^j}(y)))$
nonnegative. Put $\delta=\pi/12$ and write $y=\cos(\pi/4+u)$.
For $0<|u|\le\delta$, choose the largest $j\ge0$ with
$9^j|u|\le\delta$, so $\delta/9<9^j|u|\le\delta$.
Since $9^j\equiv1\pmod8$ and $\sin t\ge1/2$ for
$t\in[\pi/6,\pi/3]$, this single summand gives
\[
 E-G(y)\ge \frac{C_0}{4}9^{-j}(9^ju)^2
     >\frac{C_0\delta}{36}|u|\ge c_G|y-a|.
\]
For $|u|>\delta$, we have $|y-a|\ge r_0=(\sqrt3-\sqrt2)/2$.
The first summand then gives the result because $C_0r_0>c_G$.
At $u=0$ it is immediate. Oddness of $G$ proves the negative-sign case.

For every integer $L\ge1$ and $s=\pm1$ we also have the signed pullback bound
\begin{equation}\label{eq:pullback}
 |T_L(x)-sa|\ge\kappa L\dist(x,\{\xi:T_L(\xi)=sa\}),
 \qquad\kappa=4(1-a)/\pi.
\end{equation}
To see this, on $[0,\pi]$ the inequality
$|\cos t-a|\ge\kappa|t-\pi/4|$ follows from the concavity chord on
$[0,\pi/4]$ and $\sin t\ge a>\kappa$ on $[\pi/4,\pi/2]$.
On $[\pi/2,\pi]$, convexity gives $\cos t\le1-2t/\pi$, and
\[
 a-\cos t\ge 2t/\pi-(1-a)
               \ge 4(1-a)(t-\pi/4)/\pi,
\]
because $a\ge1/2$. Reflection gives the analogous assertion for $-a$.
Apply these inequalities on the half-period of $L\theta$ containing the
given angle, where $x=\cos\theta$. The unique preimage of the prescribed
target in that half-period lies in $[0,\pi]$ after division by $L$.
The inequality $|\cos\theta-\cos\phi|\le|\theta-\phi|$ proves
\eqref{eq:pullback}.

For $n\le2$ use $L=3$ and target $sa$; for $n\ge3$ use $L=9$ and
target $-sa$, with the factor $1/3$ in \eqref{eq:baseline}. Both give
\[
 E_n-sb_n(x)\ge 3c_G\kappa\dist(x,S_{n,s}),\qquad
 3c_G\kappa=\frac{3-2\sqrt2}{27\sqrt2}>\frac1{250}.
\]
The final strict inequality is equivalent to $750>527\sqrt2$, verified by
squaring positive quantities.
\end{proof}

\begin{lemma}\label{lem:lip}
For $r_n=H-v_n$, every $0\le n\le6$, every contact $\xi\in S_n$, and
every $x\in I$,
\[
 |r_n(x)-r_n(\xi)|\le 10^5|x-\xi|.
\]
This is a bound relative to each fixed contact; it asserts no global
Lipschitz regularity of $H$.
\end{lemma}
\begin{proof}
Pair the perturbation series:
\[
 H=\sum_{j\ge0}9^{-j}Q\circ T_{9^j},\qquad Q=h-h\circ T_3/3.
\]
Since $T_3'(\pm a)=3$ and $h_{\rm even}=T_4-1$, we have
$Q'(\pm a)=0$. For $x=\cos(\pi/4+u)$, $0<|u|\le\delta=\pi/12$,
choose the last scale $j$ as in Lemma~\ref{lem:gap}. Taylor's theorem and a
geometric tail bound give
\begin{align*}
 |H(x)-H(a)|
 &\le \tfrac12\norm{Q''}u^2\sum_{k=0}^{j}9^k
                  +2\norm Q\sum_{k=j+1}^\infty9^{-k}\\
 &\le \left(\tfrac{9\delta}{16}\norm{Q''}
                  +\tfrac9{4\delta}\norm Q\right)|u|.
\end{align*}
Here $|x-a|\ge|u|/2$. The same proof around $3\pi/4$ applies to $-a$.
Thus, whenever $|x\mp a|\le r_0$, a valid constant is
\[
 L_H=\tfrac{3\pi}{32}\norm{Q''}+\tfrac{54}{\pi}\norm Q<7000.
\]
For the numerical bound, $\beta<0.420$ and the Chebyshev derivative formulas
$\norm{T_k'}=k^2$, $\norm{T_k''}=k^2(k^2-1)/3$ give
\[
 \norm h<4.42,\qquad\norm{h'}<87,\qquad\norm{h''}<810,
\]
and hence $\norm Q<5.894$ and
$\norm{Q''}\le28\norm{h''}+8\norm{h'}<23376$.
Using $3<\pi<22/7$ yields the displayed bound for $L_H$.

If $L=3^k$ and $T_L(\xi)=\pm a$, iterate the functional equation $k$ times.
The derivative bound for its finite prefix is
$\norm{h'}\sum_{i=0}^{k-1}3^i=(L-1)\norm{h'}/2$.
Since $\norm{T_L'}=L^2$, this gives
\[
 |H(x)-H(\xi)|\le\left(\tfrac{L-1}{2}\norm{h'}+LL_H\right)|x-\xi|
 \quad\text{if } |x-\xi|\le r_0/L^2.
\]
For our two grids $L=3,9$, and $r_0/81>10^{-3}$.
Adding the derivative bound on $v_n$ gives a constant less than
$4\cdot87+9\cdot7000+216<10^5$ when $|x-\xi|\le10^{-3}$.
The inequality is non-strict, including $x=\xi$.

Finally $\norm H\le\tfrac32\norm h<6.63$, so $\norm{r_n}<20$.
For $|x-\xi|>10^{-3}$,
$|r_n(x)-r_n(\xi)|<40<10^5|x-\xi|$, which extends the estimate to all $I$.
\end{proof}

\section{One persistence inequality and completion}
Fix a base degree $n$ and a sign $s\in\{-1,1\}$. Choose a nearest contact
$\xi\in S_{n,s}$ and put $d=|x-\xi|$. Lemmas~\ref{lem:gap} and
\ref{lem:lip}, together with \eqref{eq:finite}, give
\begin{align*}
 s\bigl(b_n(x)+\varepsilon r_n(x)\bigr)
 &\le E_n-d/250+\varepsilon\bigl(\lambda_n+10^5d\bigr)\\
 &=E_n+\varepsilon\lambda_n-(1/250-10^5\varepsilon)d\\
 &\le E_n+\varepsilon\lambda_n
 \qquad\left(0\le\varepsilon\le\frac1{25000000}\right).
\end{align*}
Apply both signs. This proves
\[
 \norm{f_\varepsilon-(p_n^0+\varepsilon v_n)}\le E_n+\varepsilon\lambda_n.
\]
At the $n+2$ alternating equality nodes the residual is exactly
$s(\xi)(E_n+\varepsilon\lambda_n)$. Its magnitude is positive, since
$E_n>0$ and $\lambda_n>0$. The alternation theorem identifies
$P_n f_\varepsilon=p_n^0+\varepsilon v_n$. All these polynomials vanish at
zero, and \eqref{eq:induction} proves the assertion for every degree.

For nonoddness, $F$ is odd and $T_3$ interchanges $a$ and $-a$.
At both points $T_4-1=-2$, so
\[
 H(a)+H(-a)=\sum_{j\ge0}(-1/3)^j(-4)=-3.
\]
It follows that $f_\varepsilon(a)+f_\varepsilon(-a)=-3\varepsilon\ne0$,
completing Theorem~\ref{thm:main}.

\section*{Reproduction and the formalization boundary}
Run \texttt{python3 scripts/check\_certificate.py} for the independent finite
check, and \texttt{lake build} followed by \texttt{lake env lean Audit.lean}
for the formal proof.

The shared module \texttt{PolynomialCertificate.lean}
proves that exact rational coefficient certificates imply the real identities;
the concrete checks use \texttt{decide +kernel}. The series pairing proof gives
its summand explicitly to avoid expensive higher-order inference.

The Lean endpoint proves the unconditional existence of a sufficiently small
positive parameter, with only \texttt{propext}, \texttt{Classical.choice}, and
\texttt{Quot.sound}. The new generic signed persistence lemma is also checked
in Lean. The explicit constants and the range $0<\varepsilon\le4\cdot10^{-8}$
are proved in this document; they are not claimed as a quantitative Lean
corollary. The original PDF is retained for comparison. Its strict
Lipschitz inequality at coincident points is corrected here to $\le$.
\begin{thebibliography}{1}
\bibitem{sv} E. B. Saff and R. S. Varga,
\emph{Remarks on some conjectures of G. G. Lorentz},
Journal of Approximation Theory \textbf{30} (1980), 29--36.
\url{https://www.math.kent.edu/~varga/pub/paper_126.pdf}.
\end{thebibliography}
\end{document}
